\documentclass[aspectratio=169,t,13pt,usenames,dvipsnames]{beamer}
\usetheme[lily]{PaloAlto}
\usecolortheme{dolphin}
\usepackage{mathtools}
\usepackage{tikz}
\usepackage{pgfplots}
\usepackage{tikz-3dplot}
\usepackage{adjustbox}
\usetikzlibrary{quantikz}
\usepackage{xcolor}
\usepackage{bm}
\pgfplotsset{compat=newest}
\logo{%
\stackbox[c][c]{\includegraphics[height=1.15cm]{logo.pdf}\\[-3pt]{\textcopyright~2021}\\{Ron K. Cytron}}}

% \logo{\includegraphics[height=2.5cm]{logo.pdf}}
\def\ColorR#1{\textcolor{BrickRed}{#1}}
\def\ColorG#1{\textcolor{\RCthree}{#1}}
\def\ColorB#1{\textcolor{\RCtwo}{#1}}
\definecolor{links}{HTML}{2E8B57}
\hypersetup{colorlinks,linkcolor=,urlcolor=links}
\def\LinkArrow#1{%
\mbox{$\,$\href{#1}{\mbox{\vrule width 0pt height 0ex depth 0pt$\mapsto$}}}$\,$}
\def\WLOG{WLOG}
\def\True{\mbox{\texttt{True}}}
\def\False{\mbox{\texttt{False}}}
\def\Zero{\mbox{$0$}}
\def\One{\mbox{$1$}}
\def\Not#1{%
\ensuremath{\Overline{\ensuremath{#1}}}}
\def\Xor#1#2{\ensuremath{#1\oplus #2}}
\def\Nand#1#2{%
\ensuremath{\mbox{Nand}(#1,#2)}}
\def\And#1#2{%
\ensuremath{#1 \wedge #2}}
\def\Or#1#2{%
\ensuremath{#1 \vee #2}}
\def\Conj#1{%
\ensuremath{#1^{\star}}}
\def\Star#1{%
\Conj{#1}}
\def\EmptyString{\ensuremath{\lambda}}
\def\Mag#1{%
\ensuremath{|\,#1\,|}}
\def\Prob#1{%
\ensuremath{\Mag{#1}^{2}}}
\def\CNOT#1#2{%
\ensuremath{\NamedGate{CNOT}(#1,#2)}}
\def\CNOTMatrix{%
\begin{pmatrix*}[r]
1 & 0 & 0 & 0 \\
0 & 1 & 0 & 0 \\
0 & 0 & 0 & 1 \\
0 & 0 & 1 & 0\end{pmatrix*}%
}
\def\CCNOT#1#2#3{%
\ensuremath{\mbox{\NamedGate{CNOT}}(#1,#2,#3)}}
\long\def\ToDo#1{\smash{\textcolor{YellowOrange}{#1}}}
\def\ApplyGate#1#2{%
\ensuremath{#1\,#2}}
\def\Inverse#1{\ensuremath{{#1}^{-1}}}
\def\Quote#1{``#1''}
\def\NamedGate#1{%
\mbox{\bfseries{\textrm{#1}}}}
\def\Identity{\NamedGate{I}}
\def\Hadamard{\NamedGate{H}}
\def\PauliX{\NamedGate{X}}
\def\PauliY{\NamedGate{Y}}
\def\PauliZ{\NamedGate{Z}}
\def\SQBG#1#2#3#4#5{%
\ensuremath{%
#1 \begin{pmatrix*}[r] #2 & #3 \\ #4 & #5\end{pmatrix*}}}
\def\HMatrix{%
\SQBG{\RootTwo{}}{1}{1}{1}{-1}}
\def\XMatrix{%
\SQBG{\relax}{0}{1}{1}{0}}
\def\ZMatrix{%
\SQBG{\relax}{1}{0}{0}{-1}}
\def\YMatrix{%
\SQBG{\relax}{0}{-i}{i}{0}}
\def\IMatrix{%
\SQBG{\relax}{1}{0}{0}{1}}
\def\Co#1{\mathbb{C}}
\def\Implies#1#2{\ensuremath{#1 \rightarrow\ #2}}
\def\Domain#1{\ensuremath{\mathcal{#1}}}
\def\CompClass#1{\Domain{#1}}
\def\SpecialX#1{\ensuremath{#1^{\star}}}
\def\Forall#1#2{\ensuremath{\forall\ #1, #2}}
\def\Degrees#1{\ensuremath{#1^{\circ}}}
\def\RootTwo{\ensuremath{\frac{1}{\sqrt{2}}}}
\def\TwoSupOp#1#2#3{%
\ensuremath{\RootTwo{}\left(#1 #3 #2\right)}}
\def\TwoSup#1#2{%
   \TwoSupOp{\ket{#1}}{\ket{#2}}{+}}
\def\QPlus{\TwoSupOp{\QZero}{\QOne}{+}}
\def\QMinus{\TwoSupOp{\QZero}{\QOne}{-}}
\def\Set#1{\ensuremath{\left\{\,#1\,\right\}}}
\def\BigSkip{%

\bigskip

}

\def\MedSkip{%

\medskip

}
\def\SmallSkip{%

\smallskip

}
\def\SineWave#1{%
{%
\pgfmathsetmacro{\Xone}{#1}
\pgfmathsetmacro{\Xtwo}{-(#1)}
\draw (0,0) sin(1,\Xone) cos (2,0) sin(3,\Xtwo) cos (4,0);
}
}
\newcommand{\tolstrut}{%
  \vrule height\dimexpr\fontcharht\font`\A+.1ex\relax width 0pt\relax
}

\DeclareRobustCommand{\Overline}[1]{%
  \ensuremath{\overline{\raisebox{0pt}[1.07\height]{#1}}}%
}
\newenvironment{TIKZP}[1][scale=1.0]{%
\adjustbox{valign=t}\bgroup
\begin{tikzpicture}[#1]
}{%
\end{tikzpicture}
\egroup
}
\long\def\Definition#1{%
\begin{block}{Definition}
#1
\end{block}}
\long\def\Postulate#1#2{%
\begin{block}{Postulate}
\begin{description} \item[#1] #2\end{description}
\end{block}}
\long\def\TwoUnequalColumns#1#2#3#4{%
\begin{columns}%
\begin{column}{#1}#3\end{column}%
\begin{column}{#2}#4\end{column}%
\end{columns}%
}
\long\def\ThreeUnequalColumns#1#2#3#4#5#6{%
\begin{columns}%
\begin{column}{#1}#4\end{column}%
\begin{column}{#2}#5\end{column}%
\begin{column}{#3}#6\end{column}%
\end{columns}%
}
\long\def\ThreeColumns#1#2#3{%
\ThreeUnequalColumns{0.33\textwidth}{0.33\textwidth}{0.33\textwidth}{#1}{#2}{#3}%
}
\long\def\TwoColumns#1#2{%
\TwoUnequalColumns{0.5\textwidth}{0.5\textwidth}{#1}{#2}%
}
\long\def\OnlyRemark#1#2{%
\only<#1>{\Remark{#2}}}

\long\def\Remark#1{%
\begin{block}{Remark}
#1
\end{block}}
\long\def\Example#1{%
\begin{example} #1\end{example}}
\def\VV{\textit{vice versa}}
\def\QZero{\ket{\hbox to 0.7em{\hss 0\hss}}}
\def\QOne{\ket{\hbox to 0.7em{\hss 1\hss}}}
\def\QState#1{\ket{\QName{#1}}}
\def\QName#1{\ensuremath{\psi_{#1}}}
\def\TensOp{\ensuremath{\bm{\otimes}}}
\def\TensProd#1#2{\ensuremath{#1 \TensOp{} #2}}
\def\TensSupProd#1#2{\ensuremath{#1^{\bm{\otimes} #2}}}
\def\PZero{\SQB{1}{0}}
\def\POne{\SQB{0}{1}}
\def\PPlus{\ensuremath{\RootTwo{}\SQB{1}{1}}}
\def\PMinus{\ensuremath{\RootTwo{}\SQB{1}{-1}}}
\def\PHad{\ensuremath\frac{1}{\sqrt{2}}\SQB{1}{1}}
\def\PHadm{\ensuremath\frac{1}{\sqrt{2}}\SQB{1}{-1}}
\def\Intersect#1#2{\ensuremath{#1 \cap #2}}
\def\Union#1#2{\ensuremath{#1 \cup #2}}
\def\Set#1{\ensuremath{\left\{\,#1\,\right\}}}
\def\AllBits#1#2{\ensuremath{#1\in \Set{0,1}^{#2}}}
\def\DQB#1#2#3#4{%
\ensuremath{\begin{pmatrix*}[r] #1 \\ #2 \\ #3 \\ #4\end{pmatrix*}}}
\def\QQB#1#2#3#4#5#6#7#8{%
\ensuremath{\begin{pmatrix*}[r] #1 \\ #2 \\ #3 \\ #4 \\ #5 \\ #6 \\ #7 \\ #8\end{pmatrix*}}}
\def\SQB#1#2{%
\ensuremath{\begin{pmatrix*}[r] #1 \\ #2\end{pmatrix*}}}
\def\CQB#1#2{%
\ensuremath{\begin{pmatrix*}[r] #1 #2\end{pmatrix*}}}
\def\CDQB#1#2#3#4{%
\ensuremath{\begin{pmatrix*}[r] #1 #2 #3 #4\end{pmatrix*}}}
\def\TZPointDiam{2pt}
\def\TZPoint#1#2#3{%
\draw[fill=black] (#1) circle (\TZPointDiam) node[#3] {#2} ;}
\def\UnitComplexCircle{%
\draw [<->] (-1.5, 0) -- (1.5, 0)  node[right] {$\Re{}$} ;
   \draw [<->] (0,-1.5) -- (0, 1.5) node[above] {$\Im{}$} ;

   \draw (0, 0) circle (1) ;
   }
   \def\TZPEast{\TZPoint{1,0}{$+1$}{above right}}
   \def\TZPNorth{\TZPoint{0,1}{$i$}{above right}}
   \def\TZPWest{\TZPoint{-1,0}{$-1$}{above left}}
   \def\TZPSouth{\TZPoint{0,-1}{$-i$}{below right}}
\def\TZText#1#2#3{%
  \draw (#1) node [#3] {#2};
}
\def\Qiskit{\texttt{qiskit}}
\def\TwoToThe#1{\ensuremath{2^{#1}}}
\def\Exp#1{\ensuremath{e^{#1}}}
\def\ExpPhase#1{\Exp{i\,#1}}
\def\ExpNegPhase#1{\Exp{-i\,#1}}
%%
%% Not sure this is used
%%
\def\PhaseLabel#1{\phase{\fbox{#1}}}
%%
\def\ColHead#1{\begin{center}\bf #1\end{center}}
\def\Complex#1#2{\ComplexGen{#1}{#2}{+}}
\def\ComplexDiff#1#2{\ComplexGen{#1}{#2}{-}}
\def\ImaginaryM#1{\ComplexDiff{\relax}{#1}}
\def\ComplexGen#1#2#3{%
\ensuremath{#1 #3 i\relax #2}}
\def\Polar#1#2{\ensuremath{#1\hbox to 0.1pt{\hss}\Exp{i\relax#2}}}
%%
%% #1 -- lower left x
%% #2 -- lower left y
%% #3 -- upper right x
%% #4 -- upper right y
%% #5 -- num lines
\def\PFilter#1#2#3#4#5{

    \draw (#1,#2) -- (#3,#2) -- (#3,#4) -- (#1,#4) -- cycle;
    \foreach \i in {1,...,#5}
    {
        \pgfmathsetmacro{\PFilterInc}{#1 + (#3-#1) / #5 * \i}
        % \edef\PFilterInc{\pgfmathresult}
        \draw (\PFilterInc, #2) -- (\PFilterInc, #4);
    }
}
%% #1 - x
%% #2 - y
%% #3 - degrees
%% #4 - options
%% #5 - radius
\def\RadiantArrowsR#1#2#3{%
\foreach \p in {0, #1,...,360} {
   \draw[#2] (0,0)  -- (\p:#3) ;
}
}
\def\RadiantArrows#1#2{%
   \RadiantArrowsR{#1}{#2}{1}
}
\newcommand\RotateTheta[1][scale=1.0]{%
\SquareOutline{}
\begin{scope}[#1]
\fill[white] (0.5,0.5) circle (0.5);
\draw (0.5,0.5) circle (0.5);
\draw[->] (1.0,0.5) arc(0:45:0.5);
\draw (0.5,0.5) node {$\theta$};
\end{scope}
}
\def\SquareOutline{%
\path (0,0) rectangle (1,1);%
}%
\newcommand\Qif[1][scale=1.0]{%
\begin{scope}[#1]
\SquareOutline{}
\clip (0,0) rectangle(1,1);
\RotateAroundCenter{45}{%
\begin{scope}
\draw (0.15,0.15) rectangle(0.85,0.85);
\end{scope}}
\end{scope}
}
\newcommand\LightSource[1][scale=1.0]{%
\begin{scope}[#1]
\SquareOutline{}
\PhotonGun{}
\begin{scope}[shift={(0.75,0.5)}]
    \RadiantArrowsR{22.5}{color=red}{0.25}
\end{scope}
\end{scope}
}
\def\PhotonGun{%
\fill (0,0.375) rectangle (0.5, 0.625);
}
\def\HPolar{%
\begin{scope}[color=red]
\SquareOutline{}
\draw[thick] (0.5,0.4) -- ++(0.35,0)
             ++(-0.35,0.1) -- ++(0.35,0)
             ++(-0.35,0.1) -- ++(0.35,0);
% \fill (0.90,0.5) circle (0.05);
\end{scope}
}
\def\VPolar{%
\begin{scope}[color=red]
\draw[thick] (0.6,0.375) -- ++(0,.250)
             ++(0.1,-.250) -- ++(0,.250)
             ++(0.1,-.250) -- ++(0,.250);
% \fill (0.95,0.5) circle (0.05);
\end{scope}
}
\newcommand\HPolarizedLightSource[1][scale=1.0]{%
\begin{scope}[#1]
\SquareOutline{}
\PhotonGun{}
\HPolar{}
\end{scope}
}
\newcommand\VPolarizedLightSource[1][scale=1.0]{%
\begin{scope}[#1]
\SquareOutline{}
\PhotonGun{}
\VPolar{}
\end{scope}
}
\newcommand\EVBomb[1][scale=1.0]{%
\begin{scope}[#1]
\fill (0.5,0.5) circle(0.5);
\end{scope}%
}
\newcommand\EVBoom[1][scale=1.0]{%
\begin{scope}[#1]
\fill[white] (0.5,0.5) circle(0.5);
\draw[red,thick] (0.5,0.5) circle(0.5);
\clip (0,0) rectangle(1,1);
\begin{scope}
\node[scale=0.35,xshift=28pt,yshift=28pt,starburst,draw,red,fill=orange,line width=1.5pt]
{BOOM!};
\end{scope}
\end{scope}
}
\newcommand\Measurement[1][scale=1.0]{%
\begin{scope}[#1]
\draw[color=white] (0,0) rectangle(1,1);
\begin{scope}
\clip (0.25,0.25) rectangle (1,1);
\draw (0,0) circle(0.9);
\draw (0,0) circle(1.0);
\end{scope}
\draw[thick] (0,0) --(0,1);
\draw[->] (0.25,0.25) -- (0.9,0.9);
\end{scope}
}
\def\TRectangle#1#2#3#4#5#6{%
\begin{scope}[#6]
\draw [fill=#5] (#1,#2) rectangle (#3,#4);
\end{scope}
}
%% #1 -- llx
%% #2 -- lly
%% #3 -- urx
%% #4 -- ury
%% #5 -- angle
\newcommand\Mirror[1][rotate=0]{%
\clip (0,0) rectangle(1,1);
\begin{scope}[#1]
   \SquareOutline{}
   \TRectangle{0}{.4375}{1}{.5625}{Gray}{rotate=0}
\end{scope}
}
\newcommand\BeamSplitter[1][scale=1.0]{%
\begin{scope}[#1]
    \SquareOutline{}
    \TRectangle{0}{.375}{1}{.625}{SkyBlue}{rotate=0}
\end{scope}
}
\def\MZ{%
\LightSource{}
\Shift{2}{0}{\RotateAroundCenter{-45}{\BeamSplitter{}}}
\Shift{2}{-2}{\RotateAroundCenter{-45}{\Mirror{}}}
\Shift{6}{0}{\RotateAroundCenter{-45}{\Mirror{}}}
\Shift{6}{-2}{\RotateAroundCenter{-45}{\BeamSplitter{}}}
\Shift{8}{-2}{\Measurement[color=orange]{}}
\Shift{6}{-4}{\RotateAroundCenter{-90}{\Measurement[color=\RCtwo]{}}}
}
\def\Shift#1#2#3{%
\begin{scope}[shift={(#1,#2)}] #3\end{scope}
}
\def\RotateAroundCenter#1#2{%
\begin{scope}[rotate around={#1:(0.5,0.5)}]
  #2
\end{scope}
}
\def\Vskip#1{\mbox{}\vskip #1\mbox{}}
\def\Hskip#1{\mbox{}\hskip #1\mbox{}}
\author{Ron K.~Cytron}
\institute{Washington University\\Saint Louis, Missouri\\[2em] \textcopyright~2021 All rights reserved by the author}
\setbeamertemplate{footline}[frame number]{}
%\setbeamertemplate{footline}[text line]{%
%  \parbox{\linewidth}{\vspace*{-8pt}\textcopyright~2021\hfill\insertshortauthor\hfill\insertpagenumber}}
%% #1 -- lecture number
%% #2 -- lecture title
%% #3 -- lecture subtitle
%% #4 -- lecture label
\def\SetTitle#1#2#3#4{%
   \title{#2}%
   \subtitle{#3}%
   \date{\today}%
   %\begin{frame}\maketitle\end{frame}%
   \frame{\titlepage}%
   \lecture[#1]{#2: #3}{lec:#4}%
}
\def\Kaye{\href{https://dl.acm.org/doi/10.5555/1206629}{Kaye}}
\def\MikeIke{\href{https://dl.acm.org/doi/10.5555/1972505}{Nielson~and~Chuang}}

\newcounter{ProtocolDialogStep}
\newenvironment{ProtocolDialog}[3]{%
\def\Incr{\stepcounter{ProtocolDialogStep}}
\def\Ref##1{\Incr{}\ThreeUnequalColumns{#1}{#2}{#3}{\arabic{ProtocolDialogStep}. ##1}{\relax}{\relax}}%
\def\Alice##1{\Incr{}\ThreeUnequalColumns{#1}{#2}{#3}{\relax}{\arabic{ProtocolDialogStep}. ##1}{\relax}}%
\def\All##1##2##3{%
\Incr{}
\ThreeUnequalColumns{#1}{#2}{#3}{\arabic{ProtocolDialogStep}. ##1}{##2}{##3}
}
\def\Bob##1{\Incr{}\ThreeUnequalColumns{#1}{#2}{#3}{\relax}{\relax}{
\arabic{ProtocolDialogStep}. ##1}}%
\setcounter{ProtocolDialogStep}{0}%
\ThreeUnequalColumns{#1}{#2}{#3}{\textbf{Ref}}{\textbf{Alice}}{\textbf{Bob}}
}{%
}
%% #1 -- options
%% #2 -- width
%% #3 -- height
%% #4 -- signals
\newenvironment{GateBox}[4][scale=1.0]{%
\edef\Signals{#4}
\edef\Height{#3}
\edef\Width{#2}
\pgfmathsetmacro{\Max}{\Signals}
\pgfmathsetmacro{\Vsep}{\Height / \Max}
\pgfmathsetmacro{\Wlen}{\Width / 5}
\def\ComputeMyY##1{%
   \pgfmathsetmacro{\MyY}{\Height - 0.5 * \Vsep - ##1*\Vsep}}
\def\Input##1##2{%
   \pgfmathsetmacro{\Xone}{-\Wlen}
   \pgfmathsetmacro{\Xtwo}{0}
   \pgfmathsetmacro{\MyY}{\Height - 0.5 * \Vsep - ##1*\Vsep}
   \draw[->] (\Xone,\MyY) node[left] {##2} -- (\Xtwo,\MyY);
}
\def\Output##1##2{%
   \pgfmathsetmacro{\Xone}{\Width}
   \pgfmathsetmacro{\Xtwo}{\Width + \Wlen}
   \pgfmathsetmacro{\MyY}{\Height - 0.5 * \Vsep - ##1*\Vsep}
   \draw[->] (\Xone,\MyY) -- (\Xtwo,\MyY) node[right] {##2};
}
\def\BoxLabel##1{%
\node[draw, fit={(0,0) (\Width,\Height)}, inner sep=0pt, label=center:\mbox{##1}] (A) {};
}
\begin{TIKZP}[#1]
   \draw (0,0) rectangle(#2,#3);
}{%
\end{TIKZP}
}
\def\PhiTheta#1#2{%
\mbox{\ensuremath{\textcolor{\RCone}{\phi = #1}, \textcolor{\RCtwo}{\theta = #2}}}}
%
% #2 -- orcacle name
% #3 -- upper left input
% #4 -- lower left input
% #5 -- upper right output
% #6 -- lower right output
%
\newenvironment{Oracle}[6][scale=1.0]{%
\begin{GateBox}[#1]{2.5}{1.5}{2}
\BoxLabel{#2}
\Input{0}{#3}
\Input{1}{#4}
\Output{0}{#5}
\Output{1}{#6}
}{%
\end{GateBox}
}
\newenvironment{TTable}[3]{%
\def\TRow##1##2##3{%
##1 & ##2 & ##3 \\}
\begin{tabular}{c|cc}
\TRow{#1}{#2}{#3}\hline
}{%
\end{tabular}
}
%%
%%  1 name
%%  2 one x
%%  3 one f(x)
%%  4 another x
%%  5 another f(x)
\def\SBitTable#1#2#3#4#5{%
\begin{center}
    \begin{TTable}{#1}{$x$}{$f(x)$}
    \TRow{ }{#2}{#3}
    \TRow{ }{#4}{#5}
    \end{TTable}
\end{center}}
%%
%%
\newenvironment{UTable}[6]{%
\def\URow##1##2##3##4##5##6{%
##1 & ##2 & ##3  & ##4 & ##5 & ##6\\}
\begin{tabular}{c|cc|c|cc}
& \multicolumn{2}{|c|}{Inputs} &
& \multicolumn{2}{|c}{Outputs} \\
\URow{#1}{#2}{#3}{#4}{#5}{#6}\hline
}{%
\end{tabular}
}
%%
%%  1 name
%%  2 one x
%%  3 one f(x)
%%  4 another x
%%  5 another f(x)
\def\SBitTable#1#2#3#4#5{%
\begin{center}
    \begin{TTable}{#1}{$x$}{$f(x)$}
    \TRow{ }{#2}{#3}
    \TRow{ }{#4}{#5}
    \end{TTable}
\end{center}}

\newenvironment{BlochSphere}[3][scale=1.0]{
    \bgroup%
    \def\BlochSpherePhiTheta{%
    \tdplotsetrotatedcoords{0}{0}{0};O
        \path[tdplot_rotated_coords] (1.0,0,0) arc (0:115:1.0) node(phi) {};
        \draw[tdplot_rotated_coords,thick,\RCone] (0.3,0,0) arc (0:115:0.3);
        \draw[tdplot_rotated_coords,\RCone,thick,dashed] (0,0,0) -- (phi);
        \draw[\RCone] (0.4,0.4,0) node {$\phi$};
        
        \tdplotsetrotatedcoords{25}{-90}{0};
        \path[tdplot_rotated_coords] (1,0,0) arc (0:45:1) node(theta) {};
        \draw[tdplot_rotated_coords,thick,\RCtwo] (0.4,0,0) arc (0:45:0.4);
        \draw[tdplot_rotated_coords,\RCtwo,thick] (0,0,0) -- (theta);
        \draw[\RCtwo] (-0.4,0,0.3) node {$\theta$};
        
        
    }
    \def\TZPointDiam{0.7pt}%
    \def\ShowAxes{%
    \draw[dashed, tdplot_rotated_coords, black] (1,0,0) node {{\Large\bf x}} (0,1,0) node {{\Large\bf y}} (0,0,1) node {{\Large\bf z}};
    }%
    \tdplotsetmaincoords{#2}{#3}%
    \begin{tikzpicture}[tdplot_main_coords, #1]
        
        
        \shade[ball color = lightgray, opacity = 0.5] (0,0,0) circle (1cm);
        
        %%
        %% The equator
        %%
        \tdplotsetrotatedcoords{0}{0}{0};
        \draw[dashed, tdplot_rotated_coords, gray] (0,0,0) circle (1);
        
        %%
        %% Arc in the x-z plane
        %%
        \tdplotsetrotatedcoords{90}{90}{90};
        \draw[dashed, tdplot_rotated_coords, gray] (1,0,0) arc (0:180:1);
        
        %%
        %% Arc in the y-z plane
        %%
        \tdplotsetrotatedcoords{0}{90}{90};
        \draw[dashed, tdplot_rotated_coords, gray] (1,0,0) arc (0:180:1);
        
        %\tdplotsetrotatedcoords{90}{90}{90};
       % \draw[dashed, tdplot_rotated_coords, black] (1,0,0) node {{\Large\bf x}} (0,1,0) node {{\Large\bf y}} (0,0,1) node {{\Large\bf z}}; 
        
        \draw[dashed, gray] (0,0,0) -- (-1,0,0);
        \draw[dashed, gray] (0,0,0) -- (0,-1,0);
        \draw[dashed, gray] (0,0,0) -- (0,0,-1);
        
        \draw[-stealth] (0,0,0) -- (1.30,0,0);
        \draw (1.45,0,0) node {$x$};
        \draw[-stealth] (0,0,0) -- (0,1.30,0);
        \draw(0,1.45,0) node {$y$};
        \draw[-stealth] (0,0,0) -- (0,0,1.30);
        \draw (0,0,1.45) node {$z$};
        
}{%
    \end{tikzpicture}
    \tdplotsetmaincoords{0}{0}
    \egroup
}

\long\def\Exercise#1#2{%
\relax}
\def\QED{\begin{flushright}\ensuremath{\square}\end{flushright}}
\def\UGateA{%
\ensuremath{\cos(\theta/2)}}
\def\UGateB{%
\ensuremath{-\ExpPhase{\lambda}\sin(\theta/2)}}
\def\UGateC{%
\ensuremath{\ExpPhase{\phi}\sin(\theta/2)}}
\def\UGateD{%
\ensuremath{\ExpPhase{(\phi+\lambda)}\cos(\theta/2)}}
\def\UGate{%
\SQBG{\relax}{\UGateA{}}{\UGateB{}}{\UGateC{}}{\UGateD{}}}%
\def\FCirc#1{%
\begin{tikzpicture}[scale=#1]%
\SquareOutline{}%
\fill (0.5,0.5) circle(0.4);%
\end{tikzpicture}%
}%
\def\SCirc{\FCirc{0.3}}%
\newenvironment{BBKey}[0]{%
\def\STD{%
\begin{tikzpicture}[scale=0.35]%
\SquareOutline{}%
\draw[ultra thick,purple] (0,0.5) -- (1,0.5);%
\draw[ultra thick,purple] (0.5,0) -- (0.5,1);%
\end{tikzpicture}}%
\def\HDM{%
\begin{tikzpicture}[scale=0.30]%
\SquareOutline{}%
\draw[ultra thick,purple] (0,0) -- (1,1);%
\draw[ultra thick,purple] (0,1) -- (1,0);%
\end{tikzpicture}}%
\def\BBUp{%
\begin{tikzpicture}[scale=0.30]%
\SquareOutline{}%
\draw[->,thick,purple] (0.5,0) -- (0.5,1);%
\end{tikzpicture}}%
\def\BBRt{%
\begin{tikzpicture}[scale=0.30]%
\SquareOutline{}%
\draw[->,thick,purple] (0,0.5) -- (1,0.5);%
\end{tikzpicture}}%
\def\BBNe{%
\begin{tikzpicture}[scale=0.30]%
\SquareOutline{}%
\draw[->,thick,purple] (0,0) -- (1,1);%
\end{tikzpicture}}%
\def\BBSe{%
\begin{tikzpicture}[scale=0.30]%
\SquareOutline{}%
\draw[->,thick,purple] (0,1) -- (1,0);%
\end{tikzpicture}}%
\def\BBBasis{%
\begin{tabular}{c|c|c}
Basis & 0 & 1 \\
\hline
\STD{} & \BBUp{} & \BBRt{} \\
\HDM{} & \BBNe{} & \BBSe{}
\end{tabular}%
}%
}{%
\relax%
}
\def\CFCirc#1{%
  \textcolor{#1}{\FCirc{0.30}}}
\def\Agree{%
  \textcolor{green}{\FCirc{0.30}}}%
\def\DisAgree{%
   \textcolor{red}{\FCirc{0.30}}}%
\def\MetaAgree{%
   \textcolor{yellow}{\FCirc{0.30}}}%
\newenvironment{EKey}[0]{%
\def\ZZ{%
\ZRot{0}%
}%
\def\ZZPO##1{%
\ZRot{\frac{\pi}{##1}}%
}%
\def\ZZNPO##1{%
\ZRot{\ensuremath{\,-\frac{\pi}{##1}}}%
}%
\def\ZRot##1{%
\ensuremath{Z_{##1}}%
}%
}{%
}
\def\AltAlert#1#2#3{%
\alt<#1>{\alert{#2}}{#3}}
\def\RCone{Rhodamine}
\def\RCtwo{CornflowerBlue}
\def\RCthree{LimeGreen}
\def\RCfour{Tan}
\def\ColorOne#1{\textcolor{\RCone}{#1}}
\def\ColorTwo#1{\textcolor{\RCtwo}{#1}}
\def\ColorThree#1{\textcolor{\RCthree}{#1}}
\def\ColorFour#1{\textcolor{\RCfour}{#1}}
\def\NegF#1{\NegOneExp{f(#1)}}
\def\NegOneExp#1{\ensuremath{(-1)^{#1}}}
\def\RootTwoN#1{\ensuremath{\frac{1}{\sqrt{2^#1}}}}
%%
%% var, exponent
%%
\def\SumBV#1#2{\OpBV{#1}{#2}{\displaystyle\sum\limits}}
\def\OpBV#1#2#3{\ensuremath{#3_{\AllBits{#1}{#2}}}}
\def\DotP#1#2{\ensuremath{#1\boldsymbol{\cdot}#2}}
%
% for hadamard results
%
%% #1 w for H((w>)
%% #2 free variable name
%% #3 n
\newcounter{HTableCounter}
\newenvironment{HTable}[3]{%
%% ##1 v^w
%% ##2 v . w
%% ##3 result
%% ##4 basis vector 000 001 etc.
\def\TE##1##2##3##4{%
\ensuremath{##4} & \ensuremath{##1} & \ensuremath{##2} & \ensuremath{##3 1} \ensuremath{\ket{##4}}}%
\begin{tabular}{c|c|c|r}
\multicolumn{4}{c}{$\sqrt{2^{#3}}$\,\Hadamard(\ket{#1})} \\ 
$v$ & \And{v}{#1} & \DotP{v}{#1} &  \multicolumn{1}{|c}{Term} \\ \hline}{%
\end{tabular}
}



% \AtBeginLecture{\frame{\Large \insertlecture}}
%%
%% test
%%
\begin{document}
\SetTitle{23}{The Deutsch-Josza problem}{Demonstrates asymptotic quantum advantage}{23}

\begin{frame}{Overview}{What will we study here?}

\begin{itemize}[<+->]
    \item This problem and its solution are due to
    \href{https://en.wikipedia.org/wiki/David_Deutsch}{David Deutsch} and \href{https://en.wikipedia.org/wiki/Richard_Jozsa}{Richard Josza}.
    \item This problem is similar to Deutsch's problem, which we have studied.
    \item But here, $f(x)$ takes an input of $n$~bits.
    \item We care, as before, whether $f(x)$ is \emph{constant} or \emph{balanced}.
    \item This problem distinguishes the complexity classes
    \begin{description}
        \item[\href{https://complexityzoo.net/Complexity_Zoo:P}{\CompClass{P}}] problems that can be solved in polynomial time on a quantum computer
        \item[\href{https://complexityzoo.net/Complexity_Zoo:E\#eqp}{\CompClass{EQP}}] problems that can be solved \textbf{e}xactly on a \textbf{q}uantum computer in \textbf{p}olynomial time.
    \end{description}
    \item The phase-kickback trick we have just studied is the key to solving this problem efficiently on a quantum computer.
\end{itemize}
    
\end{frame}

\begin{frame}{Exact}{What's up with that?}
\begin{itemize}[<+->]
    \item There are two sources of inaccuracy in quantum computing:
    \begin{itemize}
        \item Some algorithms are inherently probabilistic, such as the Elitzur--Vaidman bomb we have studied. While you are likely to get a desired result, it may take effort to overcome the probability of obtaining the wrong answer.  And you may be unlucky to the point where it takes a long time to get the answer.
        \item Quantum computing devices are subject to noise and other interference that causes \href{https://en.wikipedia.org/wiki/Quantum_decoherence}{decoherence}.  Results from a system experiencing decoherence are unpredictable and invalid.
    \end{itemize}
    \item Simulation (emulation) provides results with no artifacts of decoherence, but those are restrictive in the number of quantum bits (\href{https://qiskit.org/documentation/tutorials/simulators/6_extended_stabilizer_tutorial.html}{currently}, 63) that can be accommodated in reasonable time.
    \item So, \emph{exact quantum polynomial} means that on an ideal quantum device, we obtain the answer exactly, with probability 100\%, in polynomial time.
\end{itemize}
\end{frame}

\begin{frame}{On the pronunciation of Josza}{From the author himself}

Thanks for your interesting question!  I generally pronounce my surname as~\Quote{joe--zuh}.
\SmallSkip{}
But this is an anglicised version---the name is Hungarian and the correct Hungarian pronunciation is  \Quote{your--zho} with \emph{zh} pronounced as the \emph{s} in \emph{pleasure} and \emph{o} as the \emph{o} in \emph{hot}.
\SmallSkip{}
I hope that gives enough of an indication, and I wish you and your students all the best for your QC course!
\BigSkip{}
Best wishes,  Richard.
\end{frame}

\begin{frame}{Deutsch--Josza}{Problem statement}
\begin{itemize}[<+->]
    \item We are given a function $f(x): \Set{0,1}^{n}\mapsto \Set{0,1}$.  It takes in a value expressed in $n$ bits and returns \True{} or \False{}.
    \item We are \href{https://en.wikipedia.org/wiki/Promise_problem}{promised} that $f(x)$ is one of the following:
    \begin{description}
        \item[constant]  \Forall{x\in \Set{0,1}^n}{f(x) = k}, where $k=0$ or $k=1$.
        \item[balanced]  For exactly half of its inputs, $f(x) = 0$ and for the other half $f(x) = 1$.
    \end{description}
\end{itemize}
\end{frame}

\begin{frame}{Sneak preview}{Of the secret sauce}
\Vskip{-3.5em}\begin{itemize}[<+->]
    \item The function $f(x)$ accepts an input \AllBits{x}{n}.  Each input is $n$~bits long, and \alert<10->{there are $2^n$ such strings} in the domain.
    \item For each such $x$, $f(x)$ evaluates to~$0$ or~$1$.  \only<10->{\alert{Then \NegF{x} is $\pm 1$.}}
    \item If we compute $\SumBV{x}{n} \alt<9->{\alert{\NegF{x}}}{f(x)}$ then we should obtain:
    \begin{description}
        \item[\alt<11->{\alert{$2^n$}}{$0$}]  We obtain this sum if \Forall{x}{f(x)=0}.
        \item[\alt<12->{\alert{$-2^{n}$}}{$2^n$}]  We obtain this sum if \Forall{x}{f(x)=1}.
        \item[\alt<13->{\alert{$0$}}{$\frac{2^n}{2}=2^{n-1}$}]  \textcolor<14->{\RCtwo}{We obtain this sum if exactly half of the function evaluations are~$1$ and the other half are~$0$.}
    \end{description}
    \item As described, we require $2^n$ function evaluations classically.  We hope to use a quantum computer to accomplish that work with a single evaluation.
    \item But first we are going to make \alert{one change} to our summation to leverage \textcolor<14->{\RCtwo}{interference} in quantum systems.
\end{itemize}
    
\end{frame}

\begin{frame}{Deutsch--Josza solution}{Compared to the Deutsch circuit}

\TwoUnequalColumns{0.55\textwidth}{0.45\textwidth}{%
\begin{center}
\Vskip{-3em}\adjustbox{valign=t, width=\textwidth}{\begin{quantikz}
\lstick{\ket{\TensSupProd{0}{n}}} & \qwbundle{\alert<1>{n}}\slice{\QState{0}} &  \gate{\alert<3>{\TensSupProd{\Hadamard}{n}}}\slice{\QState{1}} & \gate[wires=2][5em]{\mbox{$U_f$}}\gateinput{$x$}\gateoutput{$x$}\slice{\QState{2}} &\gate{\alert<3>{\TensSupProd{\Hadamard}{n}}}\slice{\QState{3}} & \meter{0/1} \\
\lstick{\alert<2>{\QZero{}}} & \gate{\alert<2>{\PauliX}} &   \gate{\alert<2>{\Hadamard}}   &  \qw\gateinput{$y$}\gateoutput{\Xor{y}{f(x)}} & \qw & \qw
\end{quantikz}} \only<1-3>{\\[1em]
Deutsch--Josza Problem, $n$ qubits}
\end{center}
}{%
\only<1-3>{%
\begin{center}
\Vskip{-3em}\adjustbox{valign=t, width=\textwidth}{\begin{quantikz}
\lstick{\QZero{}} & \qw\slice{\QState{0}} &  \gate{\Hadamard}\slice{\QState{1}} & \gate[wires=2][5em]{\mbox{$U_f$}}\gateinput{$x$}\gateoutput{$x$}\slice{\QState{2}} &\gate{\Hadamard}\slice{\QState{3}} & \meter{0/1} \\
\lstick{\QZero{}} & \gate{\PauliX} &   \gate{\Hadamard}   &  \qw\gateinput{$y$}\gateoutput{\Xor{y}{f(x)}} & \qw & \qw
\end{quantikz}} \only<1-3>{\\[1.2em]
Deutsch Problem, single qubit}
\end{center}%
}%
}
\only<1-3>{%
\BigSkip{}
\begin{itemize}
    \item<1-> The Deutsch--Josza input for $x$ is now $n$ qubits.
    \item<2-> The input for $y$ is still a single qubit, whose value will be \ket{-} going into $U_f$.
    \item<3-> We apply an $n$-way Hadamard operation on either side of $U_f$ to the top $n$~qubits~($x$).  This is realized by a single Hadamard gate applied to each qubit for $x$.
\end{itemize}
}

    
\end{frame}

\begin{frame}{Deutsch--Josza Solution}{Up to \QState{2}, prior to the final \TensSupProd{\Hadamard}{n}}
\only<4-24,32->{%
\Vskip{-4em}\TwoUnequalColumns{0.51\textwidth}{0.49\textwidth}{%
\begin{center}
\Vskip{-3em}\adjustbox{valign=t, width=\textwidth}{\begin{quantikz}
\lstick{\ket{\textcolor<8->{\RCone}{\TensSupProd{0}{n}}}} & \qwbundle{n}\slice{\alert<8>{\QState{0}}} &  \gate{\alert<5>{\TensSupProd{\Hadamard}{n}}}\slice{\alert<9-10>{\QState{1}}} & \gate[wires=2][5em]{\mbox{$U_f$}}\gateinput{$x$}\gateoutput{$x$}\slice{\alert<11>{\QState{2}}} &\gate{\alert<6,11>{\TensSupProd{\Hadamard}{n}}}\slice{\alert<12->{\QState{3}}} & \meter{0/1} \\
\lstick{\textcolor<8->{\RCtwo}{\QZero{}}} & \gate{\textcolor<8->{\RCtwo}{\PauliX}} &   \gate{\Hadamard}   &  \qw\gateinput{$y$}\gateoutput{\Xor{y}{f(x)}} & \qw & \qw
\end{quantikz}}
\end{center}
\only<11>{%
\BigSkip{}
This is the phase-kickback trick result.  But how do we compute The \alert{\TensSupProd{\Hadamard}{n}} when its input is something other than \TensSupProd{0}{n}?}
}{%
\only<8-24>{%
\Vskip{-2.5em}\begin{align*}
   \only<8-10>{ \QState{0} &= \ket{\ColorOne{\TensSupProd{0}{n}}\,\ColorTwo{1}} \\}
   \only<8-11>{%
   \visible<9->{ \QState{1} &= \TensSupProd{\Hadamard}{n+1} \QState{0} \\}
    \visible<10-> {& =  \ColorOne{\RootTwoN{n}}\TensProd{\ColorOne{\SumBV{w}{n}\ket{w}}}{\ColorTwo{\ket{-}}} \\}}
    \only<11-13>{%
    \visible<11->{\QState{2} &= \ColorOne{\RootTwoN{n}}\TensProd{\ColorOne{\SumBV{w}{n} \NegF{w}\ket{w}}}{\ColorTwo{\ket{-}}} \\}
    \visible<12->{\QState{3} &= \left(\TensProd{\TensSupProd{\Hadamard}{n}}{\Identity}\right)(\QState{2}) \\}}
    \visible<13->{%
       \only<14->{\QState{3}} & = \frac{\TensSupProd{\Hadamard}{n}
       (
       \ColorOne{\SumBV{w}{n} \NegF{w}\ket{w}}
       )}{\ColorOne{\sqrt{2^n}}} \\
       & \TensOp{} \ColorTwo{\only<13>{\Identity}\ket{-}}
    }
\end{align*}
}%
\only<32->{%
\begin{itemize}
    \item<32-> If we measure \TensSupProd{0}{n} only, then $f(x)$ is constant.
    \item<33-> Otherwise we should never see \TensSupProd{0}{n} and $f(x)$ is balanced.
    \item<34-> What if \TensSupProd{0}{n} seen sometimes?
\end{itemize}
\MedSkip{}
}
}}%
\only<4-7>{%
\begin{itemize}
    \item<4-> We have seen
    \Vskip{-3em}\[ \TensSupProd{\Hadamard}{n}\ket{\TensSupProd{0}{n}} = 
    \RootTwoN{n}\SumBV{w}{n} \ket{w} = \RootTwoN{n}\DQB{1}{\vdots}{\vdots}{1}
    \]
    \item<5-> So we know the result of these $n$ \Hadamard{} operations.
    \item<6-> But what about these?  Their input is probably not \ket{\TensSupProd{0}{n}}, so we need to analyze what happens there.
    \item<7-> We present results first and proofs later.
\end{itemize}
}%
\only<14-16>{%
\BigSkip{}
\begin{itemize}
\item<14-> How do we compute \TensSupProd{\Hadamard}{n}(\ket{w}), where \ket{w} is a basis state, but not necessarily~\ket{\TensSupProd{0}{n}}?  
\item<15-> We show how to do this here, leaving the proof for later.
\item<16-> It suffices to study basis states, as linearity applies to \ColorOne{superpositions}.
\end{itemize}
}%
\only<17-20>{%
\TwoColumns{%
\Vskip{-2.5em}\begin{center}
\begin{TIKZP}[scale=0.9]
\draw (0,0) rectangle ++(0.5,0.8) node[pos=0.5] {$v_1$};
\draw (0.5,0) rectangle ++(0.5,0.8) node[pos=0.5] {$v_2$};
\draw (1,0) rectangle ++(0.5,0.8) node[pos=0.5] {$\cdot$};
\draw (1.5,0) rectangle ++(0.5,0.8) node[pos=0.5] {$\cdot$};
\draw (2,0) rectangle ++(0.5,0.8) node[pos=0.5] {$\cdot$};
\draw (2.5,0) rectangle ++(0.5,0.8) node[pos=0.5] {$v_n$};
\draw (0,-1.2) rectangle ++(0.5,0.8) node[pos=0.5] {$w_1$};
\draw (0.5,-1.2) rectangle ++(0.5,0.8) node[pos=0.5] {$w_2$};
\draw (1,-1.2) rectangle ++(0.5,0.8) node[pos=0.5] {$\cdot$};
\draw (1.5,-1.2) rectangle ++(0.5,0.8) node[pos=0.5] {$\cdot$};
\draw (2,-1.2) rectangle ++(0.5,0.8) node[pos=0.5] {$\cdot$};
\draw (2.5,-1.2) rectangle ++(0.5,0.8) node[pos=0.5] {$w_n$};
\end{TIKZP}\end{center}
\begin{itemize}
    \item<18-> \ColorThree{This is the exclusive-or of the \emph{and} of each pair of bits.}
    \item<19-> \ColorFour{This is the sum of the product of each pair of bits.}
\end{itemize}
}{%
\SmallSkip{}
We define the operation \DotP{v}{w} on two $n$-bit vectors $v$ and $w$ as follows:
\begin{align*}
   \visible<18->{ \DotP{v}{w} & \ColorThree{= \oplus_{i=1}^{n}\ {v_i}{w_i}} \\}
        \visible<19->{ & \ColorFour{= \sum_{i=1}^{n}\ {v_i}{w_i}\ \mbox{(mod 2)}}}
\end{align*}
\only<20->{\Vskip{-1em}These are equivalent (proof later).}
}}%
\only<21-26>{%
\[
P(n):  \TensSupProd{\Hadamard}{n}\ket{w}
  = \textcolor<26>{\RCthree}{\RootTwoN{n}\SumBV{v}{n} \NegOneExp{\DotP{v}{w}} \ket{v}}
\]}%
\only<21-23>{%
\begin{itemize}
\item<22-> \TensSupProd{\Hadamard}{n}\ket{\TensSupProd{0}{n}} provides the correct result, because \Implies{w=0^n}{\Forall{v}{\DotP{v}{w}=0}}.
\item<23-> We can use this formula to compute \QState{3}.
\end{itemize}}
\only<24-31>{%
\Vskip{-2em}\begin{align*}
\only<24-26>{%
    \visible<24->{\ColorOne{\sqrt{2^n}} \QState{3} &= \TensProd{\TensSupProd{\Hadamard}{n}
       (
       \ColorOne{\SumBV{w}{n} \NegF{w}\ket{w}}
       )}{\ColorTwo{\ket{-}}} \\}
    \visible<25->{   &= \TensProd{\ColorOne{\SumBV{w}{n} \NegF{w}} \TensSupProd{\Hadamard}{n}(\ColorOne{\ket{w}})}{\ColorTwo{\ket{-}}} \\}}
    \only<26-28>{\visible<26->{
    \ColorThree{\sqrt{2^n}}\ColorOne{\sqrt{2^n}}\QState{3} &= \TensProd{\ColorOne{\SumBV{w}{n} \NegF{w}}\ColorThree{\SumBV{v}{n} \NegOneExp{\DotP{v}{w}} \ket{v}}}{\ColorTwo{\ket{-}}} \\
    }}%
    \visible<28->{%
    \only<29->{\ColorThree{\sqrt{2^n}}\ColorOne{\sqrt{2^n}}\QState{3}}&= \TensProd{\ColorThree{\SumBV{v}{n}} \ColorOne{\SumBV{w}{n}} \ColorOne{\NegF{w}} \ColorThree{\NegOneExp{\DotP{v}{w}}} \ColorThree{\ket{v}}}{\ColorTwo{\ket{-}}}
    }
\end{align*}
}%
\only<29->{%
Look at the $\ColorThree{\ket{v}}=\ket{\TensSupProd{0}{n}}$ term, where \Forall{\ColorThree{v}}{\ColorThree{\NegOneExp{\DotP{v}{w}}}=1}
    \begin{description}
      \item<30->[\Forall{w}{\ColorOne{f(w)}=0}:]  $\ColorOne{\SumBV{w}{n} 1}\ =\ {2^n}$. All other \ColorThree{\ket{v}} terms have zero amplitude.
      \item<31->[\Forall{w}{\ColorOne{f(w)}=1}:] $\ColorOne{\SumBV{w}{n} -1}\ =\ -{2^n}$. All other \ColorThree{\ket{v}} terms have zero amplitude.
      \item<31->[\ColorOne{$f(w)$} balanced:]
      \Vskip{-2em}\[
      \ColorOne{\SumBV{w}{n} \NegOneExp{f(w)}}
      = 2^{n-1} - 2^{n-1}
      = 0
      \]
    \end{description}
}
\end{frame}

\section*{Proofs}

\begin{frame}{Proof}{Hadamard on non-zero states}
    
\end{frame}

\begin{frame}{Proof}{about dot products}
    
\end{frame}
    

\end{document}
